\documentclass[10pt,a4paper]{amsart}
\usepackage[margin=19mm]{geometry}
\usepackage{amsmath,amssymb,mathtools}
\usepackage[scr=rsfs,cal=euler]{mathalfa}
\usepackage{xfrac}
\usepackage[unicode,hidelinks,bookmarks=false]{hyperref}
\newcommand{\ie}{i.e.\ }
\newcommand{\cf}{cf.\ }
\newcommand{\resp}{respectively\ }
\newcommand{\Subsection}{\subsection}
\providecommand{\nobreakdash}{\nobreak\hbox{-}}
\setlength{\emergencystretch}{2em}
\multlinegap0pt
\usepackage{amssymb}
\newtheorem{lem}{Lemma}
\newtheorem{pf}{Proof}
\newcommand{\D}{\mathbb{D}}  
\newcommand{\N}{\mathbb{N}}
\newcommand{\R}{\mathbb{R}}
\newcommand{\co}{\colon}
\newcommand{\ra}{\rightarrow}
\title{Quasiregular maps of Sierpi\'nski carpet\\ Julia sets}
\author{Sergei Merenkov, Letian Shen}
\date{Source: arXiv:2601.18109v2 (CC BY 4.0)}
\begin{document}
\maketitle

\noindent\emph{Source and licence.} Quasiregular maps of Sierpi\'nski carpet\\ Julia sets. Version arXiv:2601.18109v2 (CC BY 4.0), distributed by arXiv under the Creative Commons Attribution 4.0 International licence, http://creativecommons.org/licenses/by/4.0/, as recorded in the arXiv licence metadata for this exact version and shown on the abstract page https://arxiv.org/abs/2601.18109v2. Full licence notice: \texttt{licenses/arxiv-2601.18109-CC-BY-4.0.txt}.\par\smallskip

\noindent\textit{Editorial scope.} Counted unit: the article's lem lem:Rot (source0.tex lines 581--587). The article's complete original proof of it is delivered with it (source0.tex lines 590--602, one \texttt{proof} environment of 13 lines). No mathematics is written, repaired or re-derived: the statement and the proof are the article's own lines, in the article's own notation, and no retained line is substituted or re-flowed.  No further source-local context is needed: every cross-reference of every retained line resolves inside the two delivered spans.  Nothing was omitted from any retained span, so the package carries no omission record.  No retained line contains a citation, so the wrapper carries no bibliography and no bibliography entry.  No retained line contains a figure environment, an image-inclusion command or drawing code, so no image is delivered and the package declares no asset dependency.  Editorial formatting only, in addition: an AMS \texttt{amsart} preamble that restates the article's own declarations and macros used by the retained lines, namely the environments \texttt{lem}, \texttt{pf} on separate counters, together with the standard packages those lines need.  The retained environments print in this document as Lemma 1, Pf 1 in order of appearance, whereas the article prints the same items with the numbering of its own full document.  The complete native source of the article as distributed in the arXiv e-print for this exact version is bundled unchanged as \texttt{sources/193367/source0.tex}.
\begin{lem} \label{lem:Rot} Let $\phi\co\partial \D \ra \partial \D$ be an orientation preserving branched covering, and suppose that there   exist numbers $k,l,n\in \N$, $k\ge 2$, such that 
\begin{equation} \label{eq:basiceq}
(P_l\circ \phi)(z)= (P_n\circ \phi\circ P_k)(z) \quad \text{for $z\in \partial \D$},
\end{equation}
where $P_n(z)=z^n$ for all $n\in\N$.
Then $l=n\cdot k$  and there exists  $b\in \partial \D, m\in \N$,  such that $\phi(z)=bz^m$ for all $z\in \partial \D$.
\end{lem}

\begin{pf} By counting topological degrees of both sides, one immediately sees that 
$l=n\cdot k$. 
%Note that by definition $\phi$ is an orientation preserving branched covering of $\overline \D$ with finite degree. \\
Let $m$ be the degree of $\phi$. Then there exists a continuous function $\alpha\co \R \ra \R$ with $\alpha(t+2\pi)=\alpha(t)+2\pi m$ such that
$$ \phi(\mathrm{e}^{ \mathrm{i} t}) = \mathrm{e}^{ \mathrm{i}\alpha(t)} \quad \text{for  $t\in \R$}. $$
By  \eqref{eq:basiceq} we have 
$$ \mathrm{e}^{ \mathrm{i} l \alpha(t)} = \mathrm{e}^{ \mathrm{i} n\alpha(kt)}$$
for $t\in \R$. This implies that there exists a constant $c\in \R$ such that 
$$ \alpha(t)=\frac{n}{l}\alpha(kt)+c=\frac 1k \alpha(kt)+c  \quad \text{for  $t\in \R$}. $$
Using induction on $s\in\N\cup\{0\}$, one obtains
$$ \alpha\left(t+\frac{2\pi}{k^s}\right)=\frac 1k \alpha\left(kt+\frac{2\pi}{k^{s-1}}\right)+c=\frac 1k \alpha(kt)+\frac{2\pi}{k^s}m+c=\alpha(t)+\frac{2\pi}{k^s}m,$$  for all $t\in\R$.
Continuity of $\alpha$ then implies that $\alpha$ is a linear function and the conclusion follows.
\end{pf}

\end{document}
